%============================================================
% CH2 · P4 · PROBLEM 2 — PATH WITH MAXIMUM GOLD (LeetCode 1219)
%============================================================
\subsection{Problem 2.4.2: Path with Maximum Gold}
\problemlabel{2.4.2}{ssec:ch2-max-gold}

\problemcard{2.4.2}{Path with Maximum Gold}{LeetCode 1219}{https://leetcode.com/problems/path-with-maximum-gold/}{\exmedium\ \lc{Medium}\ (a grid with cycles: the marks are needed now, and the cost must be measured)}{Backtracking from every gold cell; mark a cell by setting it to 0 and restore it on exit; prune with the gold that is still reachable}

\subsubsection*{1. Problem statement}

A gold mine is an $m \times n$ grid; cell $(r,c)$ holds $\texttt{grid}[r][c] \ge 0$ units of gold. Starting at any cell with gold, you may repeatedly step up, down, left or right, collecting all the gold of every cell you stand on. You may never visit a cell twice and never step onto a cell with $0$ gold; you may stop anywhere. Return the maximum total gold.

\textbf{Constraints:} $1 \le m, n \le 15$; $0 \le \texttt{grid}[r][c] \le 100$; \textbf{at most 25 cells contain gold}.

\begin{table}[H]
\centering\small
\begin{tabular}{@{}lcl@{}}
\toprule
\textbf{grid} & \textbf{Output} & \textbf{A best path} \\
\midrule
\texttt{[[0,6,0],[5,8,7],[0,9,0]]} & \texttt{24} & $9 \to 8 \to 7$ \\
\texttt{[[1,0,7],[2,0,6],[3,4,5],[0,3,0],[9,0,20]]} & \texttt{28} & $1 \to 2 \to 3 \to 4 \to 5 \to 6 \to 7$ \\
\bottomrule
\end{tabular}
\end{table}

\begin{diagrambox}
\begin{tikzpicture}[x=0.75cm, y=0.75cm]
    \tikzset{gc/.style={gcell, minimum size=7.5mm, fill=answerfill!60}, zc/.style={gcell, minimum size=7.5mm, fill=lightgray, text=midgray},
             pc/.style={gc, fill=currentfill, draw=accentorange, line width=1.1pt}}
    \begin{scope}
        \foreach \v/\s [count=\i from 0] in {0/zc,6/gc,0/zc} \node[\s] at (\i,0) {\v};
        \foreach \v/\s [count=\i from 0] in {5/gc,8/pc,7/pc} \node[\s] at (\i,-1) {\v};
        \foreach \v/\s [count=\i from 0] in {0/zc,9/pc,0/zc} \node[\s] at (\i,-2) {\v};
        \draw[gpath, -{Stealth[length=2.5mm]}] (1,-1.7) -- (1,-1.3);
        \draw[gpath, -{Stealth[length=2.5mm]}] (1.3,-1) -- (1.7,-1);
        \node[font=\scriptsize\bfseries] at (1,-3.0) {Example 1: $9 + 8 + 7 = 24$};
    \end{scope}
    \begin{scope}[xshift=4.2cm, yshift=1.1cm]
        \foreach \row/\y in {{1/pc,0/zc,7/pc}/0, {2/pc,0/zc,6/pc}/-1, {3/pc,4/pc,5/pc}/-2, {0/zc,3/gc,0/zc}/-3, {9/gc,0/zc,20/gc}/-4}
            \foreach \v/\s [count=\i from 0] in \row \node[\s] at (\i,\y) {\v};
        \draw[gpath, -{Stealth[length=2.5mm]}] (0,-0.3) -- (0,-0.7);
        \draw[gpath, -{Stealth[length=2.5mm]}] (0,-1.3) -- (0,-1.7);
        \draw[gpath, -{Stealth[length=2.5mm]}] (0.3,-2) -- (0.7,-2);
        \draw[gpath, -{Stealth[length=2.5mm]}] (1.3,-2) -- (1.7,-2);
        \draw[gpath, -{Stealth[length=2.5mm]}] (2,-1.7) -- (2,-1.3);
        \draw[gpath, -{Stealth[length=2.5mm]}] (2,-0.7) -- (2,-0.3);
        \node[font=\scriptsize\bfseries] at (1,-5.0) {Example 2: 28};
    \end{scope}
    \node[note, text width=5.2cm, align=left, anchor=west] at (9.6,-1.3) {Example 2: the 20 and the 9 have no gold neighbour at all --- each is a path of one cell, worth less than the chain. The 3 below the 4 is a dead end: a path can end there, but then it misses one side of the chain.};
\end{tikzpicture}
\figcap{The two examples; grey cells hold no gold. Orange: one optimal path.}
\end{diagrambox}

\subsubsection*{2. The graph model}

\begin{modelbox}[Graph Model --- Path with Maximum Gold]
\renewcommand{\arraystretch}{1.35}
\begin{tabularx}{\linewidth}{@{}>{\bfseries}l X@{}}
Vertices & the gold cells only, $k \le 25$ of them; vertex $(r,c)$ has the value $\texttt{grid}[r][c] \in [1, 100]$ \\
Edges & undirected, between gold cells that share a side; implicit (look at the 4 neighbours); a grid graph is \textbf{bipartite} (chessboard colours) and has degree $\le 4$ \\
Weights & on the vertices, all positive \\
Question & the maximum-value \textbf{simple path}, any start, any end --- template 2 of the pattern \\
Size & $k \le 25$, so the graph is tiny; the search tree is not: up to $3\,060\,417$ nodes (measured below) \\
\end{tabularx}
\end{modelbox}

\begin{diagrambox}
\begin{tikzpicture}[x=1.2cm, y=1.1cm]
    \tikzset{gv/.append style={minimum size=7mm}}
    \node[gv] (g01) at (1,0) {6};
    \node[gv] (g10) at (0,-1) {5}; \node[gv] (g11) at (1,-1) {8}; \node[gv] (g12) at (2,-1) {7};
    \node[gv] (g21) at (1,-2) {9};
    \foreach \a/\b in {01/11, 10/11, 11/12, 11/21} \draw[ge] (g\a) -- (g\b);
    \node[font=\scriptsize\bfseries] at (1,-2.8) {Example 1: a star};
    \begin{scope}[xshift=4.6cm, yshift=0.6cm]
        \foreach \v/\n/\x/\y in {1/a/0/0, 7/b/2/0, 2/c/0/-1, 6/d/2/-1, 3/e/0/-2, 4/f/1/-2, 5/g/2/-2, 3/h/1/-3, 9/i/0/-4, 20/j/2/-4}
            \node[gv, minimum size=6.5mm, font=\scriptsize\bfseries] (n\n) at (\x,\y) {\v};
        \foreach \a/\b in {a/c, c/e, e/f, f/g, g/d, d/b, f/h} \draw[ge] (n\a) -- (n\b);
        \node[font=\scriptsize\bfseries, align=center] at (1,-4.85) {Example 2: a tree plus\\two isolated vertices};
    \end{scope}
    \node[note, text width=4.6cm, align=left, anchor=west] at (8.4,-1) {Grey cells are not vertices at all. The real graph is often a forest of small pieces; only a big \emph{compact} block of gold makes the search tree large.};
\end{tikzpicture}
\figcap{The graphs of the examples: one vertex per gold cell, one edge per pair of side-adjacent gold cells.}
\end{diagrambox}

\subsubsection*{3. Key observations}

\begin{itemize}[leftmargin=*]
    \item \textbf{The path's history matters.} ``Never visit a cell twice'' depends on the whole path; the grid has cycles (any $2 \times 2$ block of gold). So the marks are needed --- and un-marked on exit (Tool~4), because another path may use the cell.
    \item \textbf{The grid is its own \texttt{onPath} array.} Set the cell to 0 when entering: the rule ``never step on 0'' then forbids it automatically. Restore the saved value on exit.
    \item \textbf{An optimal path cannot be extended.} All values are positive, so a best path has no free gold neighbour at either end. This makes the bound ``everything still reachable from here'' strong.
    \item \textbf{How big is the search?} With $\le 25$ gold cells and degree $\le 4$ the crude bound $25 \cdot 3^{24}$ is meaningless. Measured instead: a hill-climbing search over 25-cell shapes inside a $9 \times 9$ board (moving one cell at a time, keeping the change when the number of calls grows) ended at the compact $5 \times 5$ block, with $3\,060\,417$ calls of the backtracking from all 25 starts. That is the realistic worst case: fast.
\end{itemize}

\subsubsection*{4. Approaches}

\begin{bruteforcebox}
\textbf{Idea: every simple path is a prefix of an ordering of the gold cells.} Enumerate all $k!$ orderings with \texttt{next\_permutation}; walk each one while consecutive cells are adjacent and keep the best prefix sum.

\begin{lstlisting}[style=cpp]
#include <bits/stdc++.h>
using namespace std;

// Brute force: every ordering of the gold cells; the longest prefix of side-adjacent cells is a path.
class Solution {
public:
    int getMaximumGold(vector<vector<int>>& grid) {
        vector<pair<int, int>> cells;
        for (int r = 0; r < (int)grid.size(); r++)
            for (int c = 0; c < (int)grid[0].size(); c++)
                if (grid[r][c] > 0) cells.push_back({r, c});
        int best = 0;                                  // cells is sorted: the first permutation
        do {
            int sum = 0;
            for (size_t i = 0; i < cells.size(); i++) {
                if (i > 0 && abs(cells[i].first - cells[i - 1].first) +
                             abs(cells[i].second - cells[i - 1].second) != 1) break;
                sum += grid[cells[i].first][cells[i].second];
                best = max(best, sum);
            }
        } while (next_permutation(cells.begin(), cells.end()));
        return best;
    }
};

int main() {
    vector<vector<int>> g1 = {{0, 6, 0}, {5, 8, 7}, {0, 9, 0}};
    vector<vector<int>> g2 = {{1, 0, 7}, {2, 0, 6}, {3, 4, 5}, {0, 3, 0}, {9, 0, 20}};
    cout << Solution().getMaximumGold(g1) << " " << Solution().getMaximumGold(g2) << "\n";   // 24 28
}
\end{lstlisting}

\begin{complexitybox}
\textbf{Time:} \bigO{k! \cdot k}: $5! = 120$ orderings for Example 1, $10! = 3.6 \cdot 10^6$ for Example 2, $25! \approx 1.6 \cdot 10^{25}$ in the worst case. \textbf{Space:} \bigO{k}. \textbf{Verdict:} \tle{} beyond $k \approx 11$; a perfect reference for tiny random grids, because it shares no code with the search.
\end{complexitybox}
\end{bruteforcebox}

\begin{algorithmbox}[Approach 2: backtracking with temporary zeroing]
\textbf{Idea.} Template 2 of the pattern on the grid. \texttt{dfs(r, c)} returns the most gold collectable by a path that starts at $(r,c)$ and avoids the cells currently set to 0: take this cell's gold, set it to 0, try the four neighbours, restore the gold.

\begin{lstlisting}[style=cpp]
#include <bits/stdc++.h>
using namespace std;

// Path with Maximum Gold (LeetCode 1219): backtracking from every gold cell.
// A cell on the current path is set to 0 (the rule "never step on 0" then forbids it) and restored on exit.
class Solution {
    int R, C;
    int dfs(vector<vector<int>>& g, int r, int c) {
        static const int dr[4] = {-1, 1, 0, 0}, dc[4] = {0, 0, -1, 1};
        int gold = g[r][c];
        g[r][c] = 0;                                   // enter: (r, c) is on the path
        int best = 0;
        for (int d = 0; d < 4; d++) {
            int nr = r + dr[d], nc = c + dc[d];
            if (nr >= 0 && nr < R && nc >= 0 && nc < C && g[nr][nc] > 0)
                best = max(best, dfs(g, nr, nc));
        }
        g[r][c] = gold;                                // exit: restore the gold
        return gold + best;
    }
public:
    int getMaximumGold(vector<vector<int>>& grid) {
        R = grid.size(); C = grid[0].size();
        int ans = 0;
        for (int r = 0; r < R; r++)
            for (int c = 0; c < C; c++)
                if (grid[r][c] > 0) ans = max(ans, dfs(grid, r, c));
        return ans;
    }
};

int main() {
    vector<vector<int>> g1 = {{0, 6, 0}, {5, 8, 7}, {0, 9, 0}};
    vector<vector<int>> g2 = {{1, 0, 7}, {2, 0, 6}, {3, 4, 5}, {0, 3, 0}, {9, 0, 20}};
    cout << Solution().getMaximumGold(g1) << " " << Solution().getMaximumGold(g2) << "\n";   // 24 28
}
\end{lstlisting}

\begin{complexitybox}
\textbf{Time:} one call per simple path that starts at a gold cell: 25 calls for Example 1, 66 for Example 2. Worst case found: the $5 \times 5$ block of gold, $3\,060\,417$ calls, $0.12$--$0.19$\,s here (all-equal or random values make no difference: without pruning the values do not change the search). Other shapes measured: $4 \times 6$ block 1\,603\,536 calls, $3 \times 8$ 826\,020, $2 \times 12$ 145\,100, a $1 \times 25$ line only 625. \textbf{Space:} \bigO{k} recursion depth, the grid itself is the mark array. \textbf{Verdict:} \ac{}-fast.
\end{complexitybox}
\end{algorithmbox}

\begin{invariantbox}[Why it is correct]
\textbf{Invariant.} When \texttt{dfs(r, c)} is called, the cells set to 0 by the active calls are exactly the cells of the current path before $(r,c)$; when it returns, the grid is exactly as it was at the call (every \texttt{g[r][c] = 0} is undone by the matching \texttt{g[r][c] = gold}).

\textbf{Claim.} \texttt{dfs(r, c)} returns the maximum gold of a simple path that starts at $(r,c)$ and uses only gold cells that are not on the current path. By induction on the number of available cells: the path either stops at $(r,c)$ (value \texttt{gold}) or continues to an available gold neighbour $(r', c')$ and then is a path from $(r', c')$ avoiding the current path \emph{and} $(r,c)$ --- exactly what the recursive call computes, because $(r,c)$ is 0 during it. Taking the maximum over the neighbours covers all cases.

The outer loop tries every possible first cell, so the answer is the maximum over all simple paths. \hfill$\blacksquare$
\end{invariantbox}

\begin{alternativebox}[Approach 3: memoise on (set of used cells, current cell) --- bitmask DP]
\textbf{Idea.} The future of a path depends only on which cells it has used and where it stands, so $\text{best}(S, u)$ = ``most gold still collectable from $u$ when the cells of $S$ are used'' can be memoised (\dpbook{13}{Bitmask DP}). Number the gold cells $0..k-1$; the key is $S \cdot 32 + u$ in a hash map, since $2^{25} \cdot 25$ states would not fit in an array.

\begin{lstlisting}[style=cpp]
#include <bits/stdc++.h>
using namespace std;
typedef long long ll;

// Alternative: memoise best(S, u) = max gold still collectable from u when the cells in S are used.
class Solution {
    vector<int> val;
    vector<vector<int>> adj;                           // gold cells numbered 0..k-1
    unordered_map<ll, int> memo;
    int best(ll S, int u) {
        ll key = S * 32 + u;
        auto it = memo.find(key);
        if (it != memo.end()) return it->second;
        int res = 0;
        for (int v : adj[u])
            if (!(S >> v & 1)) res = max(res, val[v] + best(S | (1LL << v), v));
        memo[key] = res;
        return res;
    }
public:
    int getMaximumGold(vector<vector<int>>& grid) {
        int R = grid.size(), C = grid[0].size();
        vector<vector<int>> id(R, vector<int>(C, -1));
        val.clear();
        for (int r = 0; r < R; r++)
            for (int c = 0; c < C; c++)
                if (grid[r][c] > 0) { id[r][c] = val.size(); val.push_back(grid[r][c]); }
        int k = val.size();
        adj.assign(k, vector<int>());
        int dr[4] = {-1, 1, 0, 0}, dc[4] = {0, 0, -1, 1};
        for (int r = 0; r < R; r++)
            for (int c = 0; c < C; c++)
                if (id[r][c] >= 0)
                    for (int d = 0; d < 4; d++) {
                        int nr = r + dr[d], nc = c + dc[d];
                        if (nr >= 0 && nr < R && nc >= 0 && nc < C && id[nr][nc] >= 0)
                            adj[id[r][c]].push_back(id[nr][nc]);
                    }
        memo.clear();
        int ans = 0;
        for (int u = 0; u < k; u++) ans = max(ans, val[u] + best(1LL << u, u));
        return ans;
    }
};

int main() {
    vector<vector<int>> g1 = {{0, 6, 0}, {5, 8, 7}, {0, 9, 0}};
    vector<vector<int>> g2 = {{1, 0, 7}, {2, 0, 6}, {3, 4, 5}, {0, 3, 0}, {9, 0, 20}};
    cout << Solution().getMaximumGold(g1) << " " << Solution().getMaximumGold(g2) << "\n";   // 24 28
}
\end{lstlisting}

\begin{complexitybox}
\textbf{Time:} one evaluation per reachable state. On the $5 \times 5$ block: $959\,085$ distinct states and $1\,187\,569$ calls --- only $2.6$ times fewer than the $3\,060\,417$ calls of Approach~2, and \emph{slower} ($0.46$--$0.75$\,s), because every call pays for a hash lookup. On a grid most pairs (set, last cell) are reached by just one path, so there is little to share. \textbf{Space:} the hash map, about a million entries. \textbf{Verdict:} \ac{}, but worse than plain backtracking here. Bitmask DP wins when the graph is dense ($K_n$: $n!$ paths but only $2^n n$ states) --- the Hamiltonian problems of Chapter~23.
\end{complexitybox}
\end{alternativebox}

\begin{tricksbox}[Approach 4: Implementation Refinements --- pruning, measured]
\textbf{Bound by the reachable gold.} At a node, flood-fill the gold cells that are still reachable from the current cell without crossing the path. If $\text{sum} + \text{reachable} \le \text{best}$, no extension can improve the answer: do not recurse. The fill costs \bigO{k} per node and is a valid optimistic bound, because any extension uses only reachable cells.

\begin{lstlisting}[style=cpp]
#include <bits/stdc++.h>
using namespace std;

// Backtracking with a bound: gold still reachable from (r, c) through free cells. If even all of it
// cannot beat the best path found so far, the branch is cut.
class Solution {
    int R, C, best;
    vector<vector<int>> g, mark;
    int stamp = 0;
    const int dr[4] = {-1, 1, 0, 0}, dc[4] = {0, 0, -1, 1};
    bool ok(int r, int c) { return r >= 0 && r < R && c >= 0 && c < C && g[r][c] > 0; }
    int reachable(int r, int c) {                      // gold reachable from the neighbours of (r, c)
        stamp++;
        int s = 0;
        vector<pair<int, int>> st(1, make_pair(r, c));
        while (!st.empty()) {
            int x = st.back().first, y = st.back().second;
            st.pop_back();
            for (int d = 0; d < 4; d++) {
                int nx = x + dr[d], ny = y + dc[d];
                if (ok(nx, ny) && mark[nx][ny] != stamp) {
                    mark[nx][ny] = stamp; s += g[nx][ny]; st.push_back(make_pair(nx, ny));
                }
            }
        }
        return s;
    }
    void dfs(int r, int c, int sum) {
        calls++;
        int gold = g[r][c];
        g[r][c] = 0; sum += gold;                      // enter
        best = max(best, sum);
        if (sum + reachable(r, c) > best)              // the bound: otherwise nothing can improve
            for (int d = 0; d < 4; d++)
                if (ok(r + dr[d], c + dc[d])) dfs(r + dr[d], c + dc[d], sum);
        g[r][c] = gold;                                // exit
    }
public:
    long long calls = 0;
    int getMaximumGold(vector<vector<int>>& grid) {
        g = grid; R = g.size(); C = g[0].size(); best = 0;
        mark.assign(R, vector<int>(C, 0));
        for (int r = 0; r < R; r++)
            for (int c = 0; c < C; c++)
                if (g[r][c] > 0) dfs(r, c, 0);
        return best;
    }
};

int main() {
    vector<vector<int>> block(5, vector<int>(5, 1));   // the worst shape for Approach 2
    vector<vector<int>> holes(4, vector<int>(6, 1));
    holes[0][1] = holes[3][4] = 0;                     // 22 cells, the best path takes 21
    for (auto* g : {&block, &holes}) {
        Solution s;
        int ans = s.getMaximumGold(*g);
        cout << ans << " " << s.calls << "\n";         // 25 65  /  21 5227
    }
}
\end{lstlisting}

\begin{center}
\small
\renewcommand{\arraystretch}{1.2}
\begin{tabular}{@{}lrrr@{}}
\toprule
\textbf{grid} & \textbf{answer} & \textbf{calls, Approach 2} & \textbf{calls, with the bound} \\
\midrule
$5 \times 5$ block, all 1 & 25 & 3\,060\,417 & 65 \\
$5 \times 5$ block, random 1..100 & 1023 & 3\,060\,417 & 65 \\
$5 \times 5$ block without its centre & 24 & 664\,288 & 63 \\
$4 \times 6$ block minus two cells & 21 & 135\,422 & 5\,227 \\
\bottomrule
\end{tabular}
\end{center}
When a path through \emph{all} gold exists, the bound stops everything as soon as one is found. The last grid is the hard case for the bound: no path takes all 22 cells, so it cannot cut until the best 21 is known, and it still saves a factor 26. The weaker bound ``everything not yet on the path'' (template 2) only helps in the first three rows; on the last grid it makes the same $135\,422$ calls as Approach~2.

\textbf{A tempting start rule.} ``An optimal path can always be chosen to start at a cell with at most 2 gold neighbours.'' On the $5 \times 5$ block it reduces the calls from $3\,060\,417$ to $614\,980$. We compared it with the full search on $260\,000$ random grids (up to $7 \times 7$, all densities, values 1--3 and 1--100, plus all-equal values) and a targeted search for the shapes that a counter-example would need, and it never lost the optimum --- but we have no proof. It is therefore \emph{not} part of the solution: the bound above is provably correct and stronger.

\textbf{Small things.} The direction arrays are \texttt{static} (or members) so they are not rebuilt per call; the grid is passed by reference --- a copy per call would cost $225$ cells at each of the 3 million nodes.
\end{tricksbox}

\subsubsection*{5. Trace}

\begin{tracebox}[{Example 1: grid = [[0,6,0],[5,8,7],[0,9,0]], directions up, down, left, right}]
\begin{center}
\begin{tikzpicture}[x=0.62cm, y=0.62cm]
    \tikzset{g2p4c/.style={gcell, minimum size=6mm, font=\scriptsize}, g2p4z/.style={g2p4c, fill=lightgray, text=midgray},
             g2p4g/.style={g2p4c, fill=answerfill!60}, g2p4p/.style={g2p4c, fill=hitfill}, g2p4u/.style={g2p4c, fill=currentfill, draw=accentorange, line width=1pt}}
    \foreach \k/\t/\a/\b/\cc/\d/\e/\vb/\vc/\vd in {
        0/{start (0,1): 6}/g2p4u/g2p4g/g2p4g/g2p4g/g2p4g/6/8/7,
        1/{6, 8: grid shows 0 0}/g2p4p/g2p4g/g2p4u/g2p4g/g2p4g/0/0/7,
        2/{6, 8, 9: leaf}/g2p4p/g2p4g/g2p4p/g2p4g/g2p4u/0/0/7,
        3/{start (1,2): 7, 8, 9}/g2p4g/g2p4g/g2p4p/g2p4p/g2p4u/6/0/0}{
        \begin{scope}[xshift=\k*3.4cm]
            \node[g2p4z] at (0,0) {0}; \node[\a] at (1,0) {\vb}; \node[g2p4z] at (2,0) {0};
            \node[\b] at (0,-1) {5}; \node[\cc] at (1,-1) {\vc}; \node[\d] at (2,-1) {\vd};
            \node[g2p4z] at (0,-2) {0}; \node[\e] at (1,-2) {9}; \node[g2p4z] at (2,-2) {0};
            \node[font=\tiny\bfseries, align=center] at (1,1.0) {\t};
        \end{scope}
    }
\end{tikzpicture}
\end{center}
{\small Yellow: on the path (already zeroed in \texttt{g}); orange: the current cell. In panel 2 the grid stores 0 at $(0,1)$ and $(1,1)$ --- that is how the search knows not to step back. Panel 4 shows the path $7, 8, 9$, the first one to reach 24.}

\begin{center}
\begin{tikzpicture}[tn/.style={rectangle, rounded corners=2pt, draw=darkgray!70, fill=white, font=\scriptsize, inner sep=2pt}]
    \node[tn] (a) at (0,0) {7 \,$\to$ 24};
    \node[tn] (b) at (0,-0.8) {8 \,$\to$ 17};
    \node[tn] (c1) at (-2.2,-1.6) {6 \,$\to$ 6};
    \node[tn, fill=currentfill, draw=accentorange] (c2) at (0,-1.6) {9 \,$\to$ 9};
    \node[tn] (c3) at (2.2,-1.6) {5 \,$\to$ 5};
    \draw[->] (a) -- (b); \draw[->] (b) -- (c1); \draw[->] (b) -- (c2); \draw[->] (b) -- (c3);
    \node[note, text width=6.2cm, anchor=west] at (3.6,-0.8) {the calls from the start $(1,2)$: each node shows the cell's gold and the value it returns, $\text{gold} + \max(\text{children})$; $7 + (8 + 9) = 24$};
\end{tikzpicture}
\end{center}

\begin{center}\small
\begin{tabular}{@{}lllr@{}}
\toprule
\textbf{start} & \textbf{calls (in order)} & \textbf{returns} & \textbf{best} \\
\midrule
$(0,1)$ gold 6 & 6, 8, 9, 5, 7 & $8 + \max(9, 5, 7) = 17$, so $6 + 17 = 23$ & 23 \\
$(1,0)$ gold 5 & 5, 8, 6, 9, 7 & $8 + \max(6, 9, 7) = 17$, so $5 + 17 = 22$ & 23 \\
$(1,1)$ gold 8 & 8, 6, 9, 5, 7 & $8 + \max(6, 9, 5, 7) = 17$ & 23 \\
$(1,2)$ gold 7 & 7, 8, 6, 9, 5 & $8 + \max(6, 9, 5) = 17$, so $7 + 17 = 24$ & \textbf{24} \\
$(2,1)$ gold 9 & 9, 8, 6, 5, 7 & $8 + \max(6, 5, 7) = 15$, so $9 + 15 = 24$ & 24 \\
\bottomrule
\end{tabular}
\end{center}
25 calls in total, 5 per start: from a leaf the search goes to 8 and then to each of the other three leaves. The answer is \textbf{24}, found from $(1,2)$ and again from $(2,1)$ --- the same path $7, 8, 9$ read in both directions.
\end{tracebox}

\begin{diagrambox}
\begin{tikzpicture}[x=0.55cm, y=0.55cm]
    \tikzset{g2p4c/.style={gcell, minimum size=5.3mm, font=\tiny}, g2p4h/.style={g2p4c, fill=answerfill!60}, g2p4p/.style={g2p4c, fill=currentfill, draw=accentorange}}
    \begin{scope}
        \foreach \x in {0,...,4} \foreach \y in {0,...,4} \node[g2p4h] at (\x,-\y) {1};
        \node[font=\scriptsize\bfseries, align=center] at (2,1.1) {$5 \times 5$ block: 3\,060\,417 calls};
    \end{scope}
    \begin{scope}[xshift=4.6cm]
        \foreach \x in {0,...,4} \foreach \y in {0,...,4} \node[g2p4c] (q\x\y) at (\x,-\y) {};
        \foreach \x/\y [count=\i] in {0/0,0/1,0/2,0/3,0/4,1/4,1/3,1/2,1/1,1/0,2/0,2/1,2/2,2/3,2/4,3/4,3/3,3/2,3/1,3/0,4/0,4/1,4/2,4/3,4/4}
            \node[g2p4p] at (\x,-\y) {\i};
        \node[font=\scriptsize\bfseries, align=center] at (2,1.1) {with the bound: the first path};
    \end{scope}
    \node[note, text width=5.6cm, anchor=west, align=left] at (13.6,-2) {The first path the search builds from the corner snakes through all 25 cells; its value equals the total, so $\text{sum} + \text{reachable} \le \text{best}$ everywhere afterwards: 65 calls in all (25 of them build this path, the rest are cut at once).};
\end{tikzpicture}
\figcap{Why the bound is so strong on the worst shape of Approach~2.}
\end{diagrambox}

\subsubsection*{6. Test suite}

\begin{testbox}
\centering\small
\begin{tabular}{@{}p{5.8cm}lp{6.2cm}@{}}
\toprule
\textbf{grid} & \textbf{Output} & \textbf{Why} \\
\midrule
\texttt{[[5]]} & 5 & one cell \\
\texttt{[[0,0],[0,0]]} & 0 & no gold at all: the loop never starts a search \\
\texttt{[[1,0,1],[0,1,0]]} & 1 & three isolated cells \\
\texttt{[[1,1],[1,1]]} & 4 & a cycle of four: the marks are essential \\
\texttt{[[10,1],[1,10]]} & 22 & the two 10s are not adjacent; every path of 4 takes all \\
$1 \times 25$ line of 1s & 25 & 625 calls: a path graph is cheap \\
$5 \times 5$ block of 1s & 25 & 3\,060\,417 calls (Approach~2), 65 with the bound \\
\bottomrule
\end{tabular}
\end{testbox}
All outputs come from running the programs. Approaches 1--4 were compared on $3000$ random grids of size up to $4 \times 4$ with at most 9 gold cells (density about 65\,\%, values 1--100): all four answers agreed every time, and the grid was always restored after the call.

\subsubsection*{7. Common mistakes}

\begin{warningbox}[Mistake 1 --- restoring the wrong value]
\begin{lstlisting}[style=snippet]
g[r][c] = 0; ... ; g[r][c] = 1;        // WRONG: the cell held gold, not 1
int gold = g[r][c]; g[r][c] = 0; ... ; g[r][c] = gold;   // RIGHT
\end{lstlisting}
The first version silently changes the mine for every later start.
\end{warningbox}

\begin{warningbox}[Mistake 2 --- a global visited array that is never cleared]
Marking cells visited ``forever'' turns the search into a plain DFS: on \texttt{[[1,1],[1,1]]} it still finds 4, but on Example 1 the first start $(0,1)$ marks all five cells and every later start finds nothing new --- the answer becomes 23 instead of 24.
\end{warningbox}

\begin{warningbox}[Mistake 3 --- starting only at the largest cell]
The best path need not contain the largest value: in Example 2 the 20 is isolated, and the answer 28 avoids it. Every gold cell must be tried as a start (or the bound must be strong enough to make that cheap).
\end{warningbox}

\begin{warningbox}[Mistake 4 --- an unproved pruning rule]
``Start only at cells with $\le 2$ gold neighbours'' passed all our random tests; ``start only at cells with \emph{exactly one} gold neighbour'' fails at once: in the $2 \times 2$ block every cell has 2 neighbours and the program would return 0. Test every rule against the unpruned search before using it.
\end{warningbox}

\subsubsection*{8. Takeaways}

\begin{itemize}[leftmargin=*]
    \item A grid with ``each cell at most once'' is a graph with cycles: marks, un-marked on exit. The grid itself can be the mark array (temporary zeroing).
    \item Do not estimate the search with $3^{24}$: \textbf{measure} it on the worst shape. Here the worst is about $3 \cdot 10^6$ calls.
    \item The best bound is ``value so far $+$ everything still reachable''. It is provable, cheap with a flood fill, and turns millions of calls into dozens.
    \item Memoising on (set, last cell) is correct but pays off only when many paths share a state --- not on sparse grids.
\end{itemize}

\begin{variantsbox}
\variant{Spell a word instead of collecting gold}{\varin{https://leetcode.com/problems/word-search/}}
  {the path must spell a given word letter by letter.}
  {the same grid search, but a cell is only entered if its letter is the next one of the word.}
  {stop at the first success; prune by letter counts (Exercise E4.2).}
\variant{Weighted edges instead of cells}{\varin{https://atcoder.jp/contests/abc317/tasks/abc317_c}}
  {a general graph of $n \le 10$ towns, roads with lengths; the longest simple route.}
  {vertices are towns, every start allowed.}
  {template 2 with edge weights instead of vertex values (Exercise E4.3).}
\variant{A path that visits every cell}{\varin{https://cses.fi/problemset/task/1625}}
  {count the 48-step paths through all 49 squares of a $7 \times 7$ grid matching a pattern.}
  {the same search with a fixed start, a fixed end and an exact length.}
  {strong dead-end pruning is needed: $46.8$ million calls without the second rule (Exercise E4.7).}
\variant{Must end on a given cell}{\varnone}
  {the path must finish at a given gold cell $t$.}
  {only nodes whose current cell is $t$ are answers.}
  {record the sum only at $t$; bound with ``reachable gold'' only if $t$ is still reachable (otherwise cut). Checked against Approach~1 restricted to orderings ending at $t$ on 1000 random grids.}
\end{variantsbox}
